\scanpage{135}

peut aussi se décrire entièrement à l'aide de
\[
\begin{aligned}
F=F_\Pi=\bigcup_{i\in\Pi}F_i=\bigl\{f\in\mathfrak P_3(S)\ \big|\ &\exists\,i\in\Pi\ \text{tel que}\\
& f\text{ soit une face}\\
&\text{pour la structure}\\
&\text{polyédrale }P_i\bigr\},
\end{aligned}\tag*{(20)}
\]
on se propose d'expliciter les conditions sur $F$, pour que $(S,F)$ provienne d'une ``structure triangulée précongruentielle'', qui sera alors néc. \uline{unique}, les $F_i$ s'interprétant alors comme les ens. des sommets des composantes connexes du graphe
\[
\Gamma_F=\{F,A,R\}\tag*{(21)}
\]
% slanted
\margin{$R$ la Relation d'inclusion $\supset$}
où on définit $A\subset\mathfrak P_2(S)$, \inserted{en termes de $F$,} comme
\[
A=\{a\in\mathfrak P_2(S)\mid\exists f\in F,\ f\supset a\}\tag*{(22)}
\]
(en fait, on aura $A=\bigcup_{i\in\Pi}A_i$) :
\[
\Pi\simeq\Pi_0(\Gamma_F).\tag*{(23)}
\]

Les axiomes sont

\deleted{(T 0)\qquad $\bigcup_{f\in F}f=S$ \unsure{(pour mémoire} -- on veut une structure triangulée !)}

(T 1)\qquad $\deleted{\unsure{\exists}}\forall\,a\in A$, existe exactement \uline{deux} $f\in F$ qui contiennent $a$,

ce qui permet de considérer $\Gamma_F$ comme un graphe. Prenant les composantes connexes

% archive sheet number: 134
